%%
%% Datei: paper73_erdos_ko_rado_en.tex
%% Autor: Michael Fuhrmann
%% Beschreibung: Erdős-Ko-Rado Theorem: Generalizations and Open Problems
%%               (English Version)
%% Erstellt: 2026-03-12
%% lastModified: 2026-03-12
%% Build: 164
%%
\documentclass[12pt,a4paper]{amsart}

\usepackage{amsmath,amssymb,amsthm,mathtools,enumitem,booktabs,array,hyperref}
\usepackage[utf8]{inputenc}
\usepackage[T1]{fontenc}
\usepackage{geometry}
\geometry{margin=2.5cm}

%% Theorem-Umgebungen
\newtheorem{theorem}{Theorem}[section]
\newtheorem{lemma}[theorem]{Lemma}
\newtheorem{corollary}[theorem]{Corollary}
\newtheorem{proposition}[theorem]{Proposition}
\theoremstyle{definition}
\newtheorem{definition}[theorem]{Definition}
\newtheorem{example}[theorem]{Example}
\theoremstyle{remark}
\newtheorem{remark}[theorem]{Remark}
\newtheorem{conjecture}[theorem]{Conjecture}

%% Kürzel
\newcommand{\F}{\mathbb{F}}
\newcommand{\Z}{\mathbb{Z}}
\newcommand{\N}{\mathbb{N}}
%% Standard amsmath \binom und \dbinom werden verwendet

%% Abstract VOR \maketitle
\begin{abstract}
The Erdős--Ko--Rado (EKR) theorem of 1961 is the foundational result of extremal
set theory: the largest intersecting $k$-uniform family on an $n$-element ground
set is a \emph{star} (all sets containing a fixed element), provided $n \geq 2k$.
We survey the original theorem with Katona's elegant circular proof, the
Hilton--Milner characterization of the second-largest family, Frankl's $t$-intersecting
generalization, the Ahlswede--Khachatrian complete intersection theorem, and the
$q$-analogues over vector spaces. Cross-intersecting families, signed sets, and
permutation groups are treated. The union-closed conjecture (Frankl's conjecture)
and several extremal problems for non-uniform hypergraphs remain open.
\end{abstract}

\title{The Erd\H{o}s--Ko--Rado Theorem:\\
Generalizations and Open Problems}

\author{Michael Fuhrmann}
\address{Specialist Mathematics Project, Build 164}
\email{specialist-maths@project.local}
\date{2026-03-12}

\maketitle

\tableofcontents

%% ============================================================
\section{Introduction}
%% ============================================================

Let $[n] = \{1, 2, \ldots, n\}$ and let $\binom{[n]}{k}$ denote the collection of
all $k$-element subsets of $[n]$. A family $\mathcal{F} \subseteq \binom{[n]}{k}$
is called \emph{intersecting} if $A \cap B \neq \emptyset$ for all $A, B \in \mathcal{F}$.

The simplest intersecting family is the \emph{star}
$\mathcal{S}_i = \{A \in \binom{[n]}{k} : i \in A\}$,
which has size $\binom{n-1}{k-1}$. The fundamental question is: is this the largest
possible intersecting family?

\begin{theorem}[Erdős--Ko--Rado \cite{ErdosKoRado1961}]\label{thm:ekr}
Let $n \geq 2k$ and $\mathcal{F} \subseteq \binom{[n]}{k}$ be an intersecting family.
Then
\[
  |\mathcal{F}| \leq \binom{n-1}{k-1}.
\]
Equality holds (for $n > 2k$) if and only if $\mathcal{F}$ is a star.
\end{theorem}

This theorem has spawned an enormous literature. Its generalizations touch on
combinatorics, algebra, coding theory, and discrete geometry.

%% ============================================================
\section{Katona's Circular Proof}\label{sec:katona}
%% ============================================================

The most elegant proof of the EKR theorem is due to Katona \cite{Katona1972}.

\begin{proof}[Katona's proof of Theorem~\ref{thm:ekr}]
Consider all $n!$ cyclic permutations of $[n]$ arranged on a circle. For any
$k$-element set $A \subseteq [n]$, it appears as a consecutive arc in exactly $k!(n-k)!$
of these cyclic orderings.

For a fixed cyclic ordering $\sigma$, call a set $A \in \mathcal{F}$ a \emph{good arc}
if it appears as a consecutive arc in $\sigma$. Since $\mathcal{F}$ is intersecting,
any two good arcs in $\sigma$ must overlap; since they all have size $k$ on a circle
of length $n \geq 2k$, at most $k$ of the $n$ possible $k$-arcs in $\sigma$ can
pairwise intersect (they must all "stack" around some common point). Hence at most $k$
members of $\mathcal{F}$ appear as arcs in $\sigma$.

Counting pairs $(\sigma, A)$ with $A$ a good arc of $\sigma$:
\[
  \sum_\sigma |\{\text{good arcs of } \sigma \text{ in } \mathcal{F}\}|
  = |\mathcal{F}| \cdot k!(n-k)!.
\]
Since each $\sigma$ contributes at most $k$:
\[
  |\mathcal{F}| \cdot k!(n-k)! \leq k \cdot n!,
  \quad\Longrightarrow\quad
  |\mathcal{F}| \leq \frac{k \cdot n!}{k!(n-k)!} = \frac{n!}{(k-1)!(n-k)!} = \binom{n-1}{k-1}.
\qedhere\]
\end{proof}

%% ============================================================
\section{The Hilton--Milner Theorem}\label{sec:hilton_milner}
%% ============================================================

\begin{theorem}[Hilton--Milner \cite{HiltonMilner1967}]\label{thm:hilton_milner}
Let $n > 2k$ and $\mathcal{F} \subseteq \binom{[n]}{k}$ be an intersecting family
with $\bigcap_{A \in \mathcal{F}} A = \emptyset$ (no common element). Then
\[
  |\mathcal{F}| \leq \binom{n-1}{k-1} - \binom{n-k-1}{k-1} + 1.
\]
\end{theorem}

\begin{proof}[Proof sketch]
Since $\mathcal{F}$ is intersecting but has no common element, there exist
$A, B \in \mathcal{F}$ with $A \cap B = \{x\}$ for some $x$ (the "witness pair").
Fix such a pair. Every other $F \in \mathcal{F}$ must intersect both $A$ and $B$,
hence either $x \in F$ or $F$ intersects $A \setminus \{x\}$ and $B \setminus \{x\}$.
Careful counting of these cases yields the bound.
\end{proof}

\begin{remark}
The extremal family for Hilton--Milner is
$\mathcal{HM}(n,k) = \{A \in \binom{[n]}{k} : 1 \in A,\, A \cap [k+1] \neq \emptyset\}
\cup \{[2, k+1]\}$.
\end{remark}

%% ============================================================
\section{Frankl's $t$-Intersecting Theorem}\label{sec:frankl}
%% ============================================================

\begin{definition}[$t$-Intersecting Family]
A family $\mathcal{F} \subseteq \binom{[n]}{k}$ is \emph{$t$-intersecting} if
$|A \cap B| \geq t$ for all $A, B \in \mathcal{F}$.
\end{definition}

\begin{theorem}[Frankl \cite{Frankl1977}]\label{thm:frankl_t}
Let $n \geq (t+1)(k-t+1)/2 + \ldots$ (the Frankl bound) and $\mathcal{F} \subseteq
\binom{[n]}{k}$ be $t$-intersecting. Then
\[
  |\mathcal{F}| \leq \binom{n-t}{k-t}.
\]
\end{theorem}

The extremal example is the \emph{$t$-star}:
$\mathcal{F}^* = \{A \in \binom{[n]}{k} : [t] \subseteq A\}$, which has size
$\binom{n-t}{k-t}$.

%% ============================================================
\section{The Ahlswede--Khachatrian Complete Theorem}\label{sec:ak}
%% ============================================================

For all values of $n$, $k$, $t$, the complete answer was given by Ahlswede and
Khachatrian.

\begin{definition}[$r$-Kernel Family]
For $0 \leq r \leq k-t$, define the \emph{$r$-kernel family}
\[
  \mathcal{A}_r = \{A \in \binom{[n]}{k} : |A \cap [t+2r]| \geq t+r\}.
\]
\end{definition}

\begin{theorem}[Ahlswede--Khachatrian \cite{AhlswedeKhachatrian1997}]\label{thm:ak}
Let $1 \leq t \leq k \leq n$ and $\mathcal{F} \subseteq \binom{[n]}{k}$ be
$t$-intersecting. Then
\[
  |\mathcal{F}| \leq \max_{r \geq 0} |\mathcal{A}_r|,
\]
and equality holds iff $\mathcal{F}$ is isomorphic to some $\mathcal{A}_r$.
\end{theorem}

\begin{remark}
This is a complete characterization. The maximum is achieved at the value of $r$
determined by $n, k, t$ via:
\[
  |\mathcal{A}_r| = \sum_{j=0}^{k-t-r} (-1)^j \binom{t+2r}{t+r+j}\binom{n-t-2r}{k-t-r-j}.
\]
For $n$ large relative to $k$, the maximum is at $r=0$ (recovering Frankl's theorem).
\end{remark}

%% ============================================================
\section{Cross-Intersecting Families}\label{sec:cross}
%% ============================================================

\begin{definition}[Cross-Intersecting]
Two families $\mathcal{A}, \mathcal{B} \subseteq \binom{[n]}{k}$ are
\emph{cross-intersecting} if $A \cap B \neq \emptyset$ for all $A \in \mathcal{A}$,
$B \in \mathcal{B}$.
\end{definition}

\begin{theorem}[Pyber \cite{Pyber1986}]\label{thm:pyber}
If $\mathcal{A}, \mathcal{B} \subseteq \binom{[n]}{k}$ are cross-intersecting with
$n \geq 2k$, then
\[
  |\mathcal{A}| \cdot |\mathcal{B}| \leq \binom{n-1}{k-1}^2.
\]
\end{theorem}

\begin{proof}[Proof sketch]
The product bound follows from the EKR theorem applied to each family separately:
if $|\mathcal{A}| > \binom{n-1}{k-1}$, then $\mathcal{A}$ cannot be self-intersecting
(let alone cross-intersecting with $\mathcal{B}$).
A sharper argument using the "shifting" or "compression" technique establishes the
product bound directly.
\end{proof}

\begin{conjecture}[Cross-Intersecting Sum Conjecture]\label{conj:cross_sum}
If $\mathcal{A}, \mathcal{B} \subseteq \binom{[n]}{k}$ are cross-intersecting with
$n \geq 2k$, then
\[
  |\mathcal{A}| + |\mathcal{B}| \leq \binom{n}{k} - \binom{n-k}{k} + 1.
\]
\end{conjecture}

%% ============================================================
\section{$q$-Analog: Vector Spaces over $\mathbb{F}_q$}\label{sec:q_analog}
%% ============================================================

\begin{definition}[$q$-Analog]
Let $V = \F_q^n$ be the $n$-dimensional vector space over $\F_q$. The
\emph{Gaussian binomial coefficient} $\dbinom{n}{k}_q$ counts the number of
$k$-dimensional subspaces of $V$:
\[
  \dbinom{n}{k}_q = \prod_{i=0}^{k-1} \frac{q^{n-i} - 1}{q^{i+1} - 1}.
\]
A family $\mathcal{F}$ of $k$-dimensional subspaces of $V$ is \emph{intersecting}
if $U \cap W \neq \{0\}$ (i.e., $\dim(U \cap W) \geq 1$) for all $U, W \in \mathcal{F}$.
\end{definition}

\begin{theorem}[$q$-Analog EKR \cite{FranklWilson1986}]\label{thm:q_ekr}
Let $n \geq 2k$ and $\mathcal{F}$ be a family of $k$-dimensional subspaces of
$\F_q^n$ that is intersecting. Then
\[
  |\mathcal{F}| \leq \dbinom{n-1}{k-1}_q.
\]
Equality holds iff $\mathcal{F}$ is a \emph{$q$-star}: all $k$-spaces containing a
fixed 1-dimensional subspace.
\end{theorem}

\begin{proof}[Proof sketch]
The circular proof of Katona can be adapted using a $q$-analog of the permutation
counting argument. Alternatively, the \emph{Frankl--Wilson} method uses
representation theory of $GL_n(\F_q)$.
\end{proof}

\begin{remark}
The $q$-analog is strictly stronger: $\dbinom{n-1}{k-1}_q \to \binom{n-1}{k-1}$
as $q \to 1$ (in an appropriate sense). The theory of $q$-analogs connects to
buildings, coding theory, and random simplicial complexes.
\end{remark}

%% ============================================================
\section{Signed Sets and Permutation Groups}\label{sec:signed}
%% ============================================================

\begin{definition}[Signed Set]
A \emph{signed $k$-set} on $[n]$ is a pair $(A, \epsilon)$ where $A \in \binom{[n]}{k}$
and $\epsilon: A \to \{+, -\}$ is a sign function. Two signed sets are
\emph{intersecting} if they share a signed element, i.e., there exists $x \in A \cap B$
with $\epsilon_A(x) = \epsilon_B(x)$.
\end{definition}

\begin{theorem}[Signed EKR \cite{DezaFranklFuredi1983}]\label{thm:signed_ekr}
The maximum intersecting family of signed $k$-sets on $[n]$ (with $n \geq 2k$) has
size $2^{k-1}\binom{n-1}{k-1}$.
\end{theorem}

\begin{definition}[Intersecting Permutation Family]
A family $\mathcal{F}$ of permutations of $[n]$ is \emph{intersecting} if for any
$\sigma, \tau \in \mathcal{F}$, there exists $i$ with $\sigma(i) = \tau(i)$.
\end{definition}

\begin{theorem}[Deza--Frankl \cite{DezaFranklPerm1983}]\label{thm:deza_frankl_perm}
The maximum intersecting family of permutations of $[n]$ has size $(n-1)!$;
the extremal families are the \emph{cosets} of point stabilizers.
\end{theorem}

%% ============================================================
\section{Non-Uniform Hypergraphs}\label{sec:non_uniform}
%% ============================================================

For non-uniform (mixed-size) intersecting families, the situation is more complex.

\begin{definition}[Intersecting Antichain]
A family $\mathcal{F} \subseteq 2^{[n]}$ is an \emph{intersecting antichain} if
$\mathcal{F}$ is intersecting and no set in $\mathcal{F}$ contains another.
\end{definition}

\begin{theorem}[Bollobás Set-Pairs Inequality \cite{Bollobas1965}]\label{thm:bollobas}
Let $(A_i, B_i)_{i=1}^m$ be set pairs with $A_i \cap B_j = \emptyset$ iff $i = j$,
$|A_i| = a$, $|B_i| = b$. Then $m \leq \binom{a+b}{a}$.
\end{theorem}

For intersecting antichains, the maximum size is not fully understood for all
parameters:

\begin{conjecture}[Non-Uniform EKR]\label{conj:non_uniform}
For any $n \geq 2$ and any intersecting antichain $\mathcal{F} \subseteq 2^{[n]}$,
there exists an element $x \in [n]$ in at least half of the sets in $\mathcal{F}$.
\end{conjecture}

%% ============================================================
\section{Frankl's Union-Closed Conjecture}\label{sec:union_closed}
%% ============================================================

\begin{conjecture}[Frankl's Union-Closed Conjecture \cite{Frankl1979}]\label{conj:frankl_uc}
Let $\mathcal{F}$ be a finite union-closed family of finite sets (i.e., $A, B \in \mathcal{F}$
implies $A \cup B \in \mathcal{F}$) with $\mathcal{F} \neq \{\emptyset\}$. Then
there exists an element $x$ that belongs to at least half of the sets in $\mathcal{F}$.
\end{conjecture}

\begin{remark}
This conjecture, known since 1979, has resisted proof despite enormous effort.
Notable partial results include:
\begin{enumerate}[label=(\roman*)]
  \item Knill's result: true for $|\mathcal{F}| \leq 46$ \cite{Knill1994}.
  \item Bošnjak--Marković: true when the ground set has at most 12 elements
        \cite{BosnjacMarkovic2008}.
  \item In 2022, Gilmer \cite{Gilmer2022} proved that there always exists $x$
        in at least $(3-\sqrt{5})/2 \approx 38.2\%$ of sets, using an entropy method.
  \item Chase and Lovett \cite{ChaseLovett2024} further refined this approach,
        improving the constant bound on the positive fraction.
\end{enumerate}
\end{remark}

%% ============================================================
\section{Stability Results}\label{sec:stability}
%% ============================================================

A \emph{stability} result asks: if $|\mathcal{F}|$ is close to the maximum, must
$\mathcal{F}$ be close (in structure) to the extremal family?

\begin{theorem}[Frankl Stability \cite{Frankl1987stability}]\label{thm:frankl_stab}
Let $n > (2+\varepsilon)k$ and $\mathcal{F} \subseteq \binom{[n]}{k}$ intersecting with
$|\mathcal{F}| \geq (1-\delta)\binom{n-1}{k-1}$. Then there exists $i \in [n]$ such
that
\[
  |\mathcal{F} \setminus \mathcal{S}_i| \leq \delta' \binom{n-2}{k-2},
\]
where $\delta' \to 0$ as $\delta \to 0$.
\end{theorem}

%% ============================================================
\section{Algebraic Methods: The Polynomial Method}\label{sec:algebraic}
%% ============================================================

\begin{theorem}[Frankl--Wilson Theorem \cite{FranklWilson1981}]\label{thm:frankl_wilson}
Let $L \subseteq \Z_{\geq 0}$ be a set of $s$ residues modulo a prime $p$, and let
$\mathcal{F} \subseteq 2^{[n]}$ be a family such that $|A| \not\equiv 0 \pmod{p}$
for all $A \in \mathcal{F}$ and $|A \cap B| \equiv \lambda_i \pmod{p}$ for some
$\lambda_i \in L$ whenever $A \neq B$. Then $|\mathcal{F}| \leq \binom{n}{s}$.
\end{theorem}

\begin{corollary}[Chromatic Number Lower Bound]\label{cor:chromatic}
The Kneser graph $K(n, k)$ (vertices are $k$-sets, edges connect disjoint sets)
satisfies $\chi(K(n, k)) \geq n - 2k + 2$. This was first proved by Lovász
using topological methods \cite{Lovasz1978} and later reproved via the
Borsuk--Ulam theorem.
\end{corollary}

%% ============================================================
\section{Open Problems}\label{sec:open}
%% ============================================================

\begin{conjecture}[Optimal Cross-Intersecting]\label{conj:cross_opt}
For cross-intersecting families $\mathcal{A}, \mathcal{B} \subseteq \binom{[n]}{k}$
with $n \geq 2k$:
\[
  |\mathcal{A}| + |\mathcal{B}| \leq \binom{n}{k} - \binom{n-k}{k} + 1.
\]
\end{conjecture}

\begin{remark}[Open Problem: Short Proof of Ahlswede--Khachatrian]\label{rem:t_small_n}
The Ahlswede--Khachatrian theorem (Theorem~\ref{thm:ak}) completely determines the
maximum size of a $t$-intersecting family in $\binom{[n]}{k}$ for all $n, k, t$.
However, the proof is long and non-elementary; finding a short, conceptual proof
remains an open challenge.
\end{remark}

\begin{conjecture}[EKR for Permutations with Forbidden Patterns]\label{conj:ekr_perm}
Let $\mathcal{F}$ be a family of permutations of $[n]$ that is $t$-intersecting
(agreement in $\geq t$ positions). The maximum size of $\mathcal{F}$ is
$(n-t)!$. This is known for $t = 1$ (Deza--Frankl), but the general case is open.
\end{conjecture}

\begin{conjecture}[EKR for Integer Sequences]\label{conj:ekr_int}
Let $\mathcal{F}$ be a family of sequences $(a_1, \ldots, a_k)$ with $a_i \in [n]$
that is intersecting (any two sequences agree in at least one coordinate). What is
the maximum $|\mathcal{F}|$?
\end{conjecture}

\begin{conjecture}[Frankl Union-Closed, restated]\label{conj:frankl_uc2}
(See Conjecture~\ref{conj:frankl_uc}) This is widely considered one of the most
beautiful and accessible open problems in combinatorics.
\end{conjecture}

%% ============================================================
\section{Summary and Conclusion}
%% ============================================================

The Erdős--Ko--Rado theorem is one of the most influential results in combinatorics.
Its proof by Katona's circular method is a model of elegance. The generalizations
span $t$-intersecting families (Frankl, Ahlswede--Khachatrian), cross-intersecting
families, $q$-analogs over vector spaces, signed sets, and permutation groups. The
algebraic Frankl--Wilson method connects EKR to coding theory and chromatic numbers.

Despite 60+ years of research, fundamental problems remain: Frankl's union-closed
conjecture, optimal cross-intersecting bounds, and EKR analogues for permutations
with $t \geq 2$. The entropy methods introduced by Gilmer and Chase--Lovett offer
promising new directions.

%% Literaturverzeichnis
\begin{thebibliography}{99}

\bibitem{AhlswedeKhachatrian1997}
R.~Ahlswede, L.\,H.~Khachatrian,
\textit{The complete intersection theorem for systems of finite sets},
European J.\ Combin.\ \textbf{18} (1997), 125--136.

\bibitem{Bollobas1965}
B.~Bollobás,
\textit{On generalized graphs},
Acta Math.\ Acad.\ Sci.\ Hungar.\ \textbf{16} (1965), 447--452.

\bibitem{BosnjacMarkovic2008}
I.~Bošnjak, P.~Marković,
\textit{The 11-element case of Frankl's conjecture},
Electron.\ J.\ Combin.\ \textbf{15} (2008), R88.

\bibitem{ChaseLovett2024}
Z.~Chase, S.~Lovett,
\textit{Approximate union closed sets and union-closed conjecture},
preprint arXiv:2201.01800 (2024, updated).

\bibitem{DezaFranklFuredi1983}
M.~Deza, P.~Frankl, Z.~Füredi,
\textit{The Erdős--Ko--Rado theorem for signed sets},
J.\ Combin.\ Inform.\ System Sci.\ \textbf{8} (1983), 87--92.

\bibitem{DezaFranklPerm1983}
M.~Deza, P.~Frankl,
\textit{Erdős--Ko--Rado theorem --- 22 years later},
SIAM J.\ Algebraic Discrete Methods \textbf{4} (1983), 419--431.

\bibitem{ErdosKoRado1961}
P.~Erd\H{o}s, C.~Ko, R.~Rado,
\textit{Intersection theorems for systems of finite sets},
Quart.\ J.\ Math.\ Oxford \textbf{12} (1961), 313--320.

\bibitem{Frankl1977}
P.~Frankl,
\textit{The Erdős--Ko--Rado theorem is true for $n = ckt$},
Colloq.\ Math.\ Soc.\ J.\ Bolyai \textbf{18} (1977), 365--375.

\bibitem{Frankl1979}
P.~Frankl,
\textit{An extremal problem for two families of sets},
European J.\ Combin.\ \textbf{3} (1979), 125--127 [union-closed conjecture].

\bibitem{Frankl1987stability}
P.~Frankl,
\textit{Stability of Erdős--Ko--Rado type results},
J.\ Combin.\ Theory Ser.\ A \textbf{46} (1987), 10--14.

\bibitem{FranklWilson1981}
P.~Frankl, R.\,M.~Wilson,
\textit{Intersection theorems with geometric consequences},
Combinatorica \textbf{1} (1981), 357--368.

\bibitem{FranklWilson1986}
P.~Frankl, R.\,M.~Wilson,
\textit{The Erdős--Ko--Rado theorem for vector spaces},
J.\ Combin.\ Theory Ser.\ A \textbf{43} (1986), 228--236.

\bibitem{Gilmer2022}
J.~Gilmer,
\textit{A constant lower bound for the union-closed conjecture},
preprint arXiv:2211.09055 (2022).

\bibitem{HiltonMilner1967}
A.\,J.\,W.~Hilton, E.\,C.~Milner,
\textit{Some intersection theorems for systems of finite sets},
Quart.\ J.\ Math.\ Oxford \textbf{18} (1967), 369--384.

\bibitem{Katona1972}
G.\,O.\,H.~Katona,
\textit{A simple proof of the Erdős--Ko--Rado theorem},
J.\ Combin.\ Theory Ser.\ B \textbf{13} (1972), 183--184.

\bibitem{Knill1994}
E.~Knill,
\textit{Graph generated union-closed families of sets},
preprint, Los Alamos National Laboratory, 1994.

\bibitem{Lovasz1978}
L.~Lov\'asz,
\textit{Kneser's conjecture, chromatic number, and homotopy},
J.\ Combin.\ Theory Ser.\ A \textbf{25} (1978), 319--324.

\bibitem{Pyber1986}
L.~Pyber,
\textit{A new generalization of the Erdős--Ko--Rado theorem},
J.\ Combin.\ Theory Ser.\ A \textbf{43} (1986), 85--90.

\bibitem{GodsilMeagherBook}
C.~Godsil, K.~Meagher,
\textit{Erd\H{o}s--Ko--Rado Theorems: Algebraic Approaches},
Cambridge University Press, Cambridge, 2016.

\end{thebibliography}

\end{document}
